% !TEX root = ../../../main.tex
% Appendix entry: preprocessing.wavelet_transform -- "Wavelet transform" (Appendix A).
\entryfield{Algorithm}
Stage~1 of the three-block \texttt{dehshibi\_spikes} template (defined in
\code{webui/server/templates.py}, \code{CANONICAL}):
\begin{center}
\blockref{preprocessing.wavelet_transform} $\to$ \blockref{detection.wavelet_summation}
$\to$ \blockref{detection.summation_threshold}.
\end{center}
\blockref{detection.dehshibi_spikes} is the same pipeline as one deprecated block.

The block ports the first half of \code{iApplyWavelet.m} from the authors' own MATLAB
(Zenodo, doi:10.5281/zenodo.3997031), not the printed Eqs.~(2)--(3) of
\citet{dehshibi2021electrical}. \code{_run} in \code{Adapters/preprocessing_wavelet_transform.py}
calls \code{transform} in \code{Working/Detection/analysis/dehshibi_authors.py}:
\begin{enumerate}
\item \emph{Cut the span into chunks} (\code{chunks_for}). By default, consecutive windows of
  $w = \mathrm{round}(\texttt{window\_s}\cdot f_s)$ samples (\code{fixed_windows}); a tail
  shorter than $w/2$ joins the window before it, so the last window can be up to $1.5w$
  long. With \emph{slice\_by\_state}, \code{split_signal} ports the authors'
  \code{iSplitSignal}: two state levels are the modes of the lower and upper halves of a
  100-bin amplitude histogram (\code{state_levels}, MATLAB's \code{statelevels}); a
  transition counts only when the signal passes from within 2\,\% of one level to within
  2\,\% of the other (IEEE Std~181, as MATLAB's \code{pulsesep}); the span is cut at the
  rising crossing of each pulse that is followed by another. Either way every sample lies in
  exactly one chunk.
\item \emph{One frequency grid for the whole span} (\code{morse_scales}): the grid MATLAB's
  default \code{cwt} builds for a signal of $N = \min(w, n)$ samples (or $w$ with
  \emph{slice\_by\_state}): analytic Morse wavelet with $\gamma = 3$, $\beta = 20$
  ($\beta\gamma = 60$), ten voices per octave, from the scale whose filter has fallen to half
  its peak at Nyquist up to the scale whose wavelet is two time-standard-deviations wide over
  $N$ samples. A 3000-sample window gives 87 rows.
\item \emph{Transform each chunk of at least 4 samples} (\code{_modulus_rows}): extend it by
  reflection (half its length on each side up to 100{,}000 samples, otherwise
  $\lceil\log_2 n\rceil$), take the FFT, multiply by the Morse filter at each scale (positive
  frequencies only, L1-normalised, peak 2), invert, keep the modulus $|W(s,\tau)|$. One scale
  at a time; the whole complex matrix never exists.
\item \emph{Scale each row of the chunk} (\code{_scale_row}, the authors' Eq.~(3)):
  \[ G(s,\tau) = 1 + \mathrm{fix}\!\left(\eta\,\frac{|W(s,\tau)| - \min_\tau |W(s,\cdot)|}
     {\max_\tau\bigl(|W(s,\cdot)| - \min_\tau |W(s,\cdot)|\bigr)}\right), \]
  with $\eta = 240$, so values are whole numbers from 1 to 241; a row with no range is all
  ones. The minimum and maximum are taken over the chunk, not the span.
\item Write the rows into a float32 image, \emph{reversed} so that row~0 is the lowest
  frequency.
\end{enumerate}

\emph{How it connects.} \blockref{detection.wavelet_summation} sums a band of this image's
rows into $\Omega(\tau)$. The authors sum the rows whose frequency is at most a quarter of the
frequency range (\code{low_band_rows}); because an \texttt{Encoding} carries no frequency axis,
this block reports that band as a fraction of the image height, in its readout and in
\code{meta["authors_band"]["row_to"]}, and the template passes $\texttt{row\_to} = 0.7586$
(66 of 87 rows) to the summation.

\entryfield{Assumptions}
\begin{itemize}
\item \emph{Units.} The block is not unit-aware, and needs not be: each row is rescaled by its
  own minimum and maximum, and the Morse filter has no response at zero frequency, so the
  image does not change if the signal is multiplied by a constant or shifted by one.
\item \emph{Sampling rate.} \emph{window\_s} is in seconds and converted with $f_s$. The
  number of rows depends on the window's length in samples, so the same 3000\,s window
  makes more rows at 10\,Hz than at 1\,Hz (\unverified{about 120}). A window needing more
  than 128 rows is refused in words (\code{MAX_ROWS}); on the default window that happens at
  \unverified{roughly 20\,Hz and above}.
\item \emph{Missing values.} The signal is read with \code{np.asarray(x, dtype=float)} and
  NaNs are kept. One NaN in a window makes its FFT, and so every value of that window's image,
  NaN. A chunk shorter than 4 samples is left NaN too. Such windows are counted in
  \code{meta["n_chunks_skipped"]}; downstream they carry ``no evidence'', not zero.
\item \emph{Span length.} With fixed windows, a span shorter than one window is transformed
  as one chunk on its own grid.
\end{itemize}

\entryfield{Outputs}
An \texttt{Encoding} of kind \texttt{image}: a float32 array of shape (rows, $n$), one column
per sample of the span, rows low to high frequency. \code{meta} holds \code{n_chunks},
\code{n_chunks_skipped}, \code{n_rows}, \code{chunks} (half-open, span-relative),
\code{freqs_hz} (one per row, ascending) and \code{authors_band} (\code{rows}, \code{row_to}).
The readout beside the settings (\code{_derive}) shows the number of windows, of frequency
rows (warned above 128), and the authors' band as a \code{row_to} fraction.

\entryfield{Cost}
FFT-bound: about $O(R\,n\log w)$ for $R$ rows. The block declares \code{max_span_samples} =
320{,}000: the image costs $128 \times 4 = 512$ bytes per sample at the row ceiling, plus about
112 bytes of transient per sample, 624 in all, against the 200\,MB budget the gramian blocks
use. That is 88.9\,h at 1\,Hz but 8.9\,h at 10\,Hz; \code{Working.execution.execute_recipe}
refuses a longer span before allocating. The time estimate (\code{_estimate}, cost model
exponent 1.1) answers nothing until \code{Working.block_cost.calibrate()} has been run on the
machine.

\entryfield{Failure modes}
\begin{itemize}
\item \emph{Every window looks busy.} The per-row, per-window rescaling stretches whatever a
  window holds to 1--241, so a flat, noisy window paints as brightly as one with events.
  This is the method's, and it carries into the extrema of $\Omega$ downstream.
\item \emph{A step at a window boundary.} Neighbouring windows are scaled independently, so
  the image can change level abruptly at a window edge.
\item \emph{Mismatched windows across the template.} \emph{window\_s} and
  \emph{slice\_by\_state} also exist on \blockref{detection.summation_threshold}, which cuts the
  score into its own windows. The template leaves both at their defaults; change one and not
  the other and the threshold's windows no longer match the transform's normalisation.
\item \emph{A stale band.} The template's \code{row_to} of 0.7586 is right only for a
  3000-sample window. At another rate or window, the summation sums a different band from the
  authors' unless you copy the fraction this block prints.
\item \emph{State slicing on the nominal grid.} With \emph{slice\_by\_state}, a piece longer
  than \emph{window\_s} is transformed on the nominal window's grid, where MATLAB's
  \code{cwt} would add lower frequencies (fixup-J report, \S10).
\item \emph{A lost window.} A single NaN blanks a whole window (3000\,s by default).
\end{itemize}

\entryfield{Validation}
\code{tests/test_dehshibi_authors.py} holds the core to what the authors' unmodified MATLAB
returned in MATLAB R2025b on three synthetic recordings with known events
(\code{tests/fixtures/dehshibi/reference.json}, made by \code{scripts/fixup_j_synthetic.py}):
the frequency grid to $10^{-9}$, the row sums of the modulus to $10^{-6}$, and the band sum
$\Omega$ to within one count. \code{tests/test_adapter_wavelet_transform.py} checks the block
against the same fixture: one column per sample, rows low to high, values whole numbers
1--241, and the authors' band. A further test checks that the reflection padding stops an
end-to-start jump from dominating the image. How well the whole template detects spikes is in
\blockentry{detection.summation_threshold}.

\entryfield{References}
\citet{dehshibi2021electrical}; the authors' code, Zenodo doi:10.5281/zenodo.3997031. The
clause-by-clause comparison of the paper, the code and the blocks is
\code{docs/prompts/fixup/reports/J-dehshibi-vs-the-paper.md}, \S4.

\entryfield{Tests}
\code{tests/test_adapter_wavelet_transform.py} (types, defaults, the image against MATLAB,
window-by-window agreement with the core, NaN for a span too short to transform, the span cap,
refusal of a window needing too many rows); \code{tests/test_dehshibi_authors.py} (the core);
\code{tests/test_template_dehshibi.py} (the template end to end).
